% MechRef Audit — standalone \input{} file
% Requires: booktabs, array, multirow, xcolor, cleveref packages in parent document

% Status symbols — \providecommand so these are harmless if
% mechval_scorecard.tex is loaded first (which uses the same names).
\providecommand{\cmark}{\textcolor{green!60!black}{\checkmark}}%
\providecommand{\pmark}{\textcolor{orange!80!black}{$\sim$}}%
\providecommand{\xmark}{\textcolor{red!70!black}{$\times$}}%

\section{MechRef Audit: Reference Transport}
\label{app:mechref}

MechRef \citep{tower2026mechreference} asks: when does a mechanistic term
pick out the same thing across systems? The framework defines a cumulative
transport hierarchy---behavioral, role, subspace, structural, object---where
higher levels impose stricter identity conditions. A term that transports at
the role level (same functional role, possibly different components) but not
at the object level (different specific parameters) has role-level reference
but not object-level reference. This appendix audits the paper's key
mechanistic terms for their transport properties.

\begin{table*}[t]
\centering
\caption{MechRef transport audit for four key terms. Each cell indicates
whether the term successfully refers to the same entity across the specified
comparison. \cmark~= transports, \pmark~= partial/conditional,
\xmark~= fails. The cumulative hierarchy means failure at level $\ell$
implies failure at all higher levels.}
\label{tab:mechref}
\small
\setlength{\tabcolsep}{4pt}
\begin{tabular}{@{}lp{2.4cm}ccccc@{}}
\toprule
\textbf{Term} & \textbf{Comparison} & \textbf{Behav.} & \textbf{Role} & \textbf{Subspace} & \textbf{Structural} & \textbf{Object} \\
\midrule
\multirow{3}{2.6cm}{``Grassmannian causal variable''}
  & Across seeds (same op)    & \cmark & \cmark & \xmark & \pmark & \xmark \\
  & Across operations         & \cmark & \cmark & \xmark & \xmark & \xmark \\
  & Across architectures      & ?      & ?      & ?      & ?      & ? \\
\midrule
\multirow{3}{2.6cm}{``Grokking''}
  & Always $\to$ always ops   & \cmark & \cmark & --- & --- & --- \\
  & Stochastic ops (diff.\ seeds) & \xmark & \xmark & --- & --- & --- \\
  & Transformer $\to$ DMF+ReLU    & \cmark & \pmark & --- & --- & --- \\
\midrule
\multirow{2}{2.6cm}{``Structured pi-SAE variable''}
  & Grokking $\to$ GPT-2      & \cmark & \cmark & \xmark & \xmark & \xmark \\
  & Across GPT-2 tasks        & \cmark & \cmark & \xmark & \xmark & \xmark \\
\midrule
\multirow{2}{2.6cm}{``Equivariance''}
  & Across symmetry groups    & \cmark & \pmark & --- & --- & --- \\
  & Algebraic $\to$ NLP tasks & \xmark & \xmark & --- & --- & --- \\
\bottomrule
\end{tabular}
\end{table*}

\paragraph{Grassmannian causal variable: object-level transport fails.}
The subspace degeneracy result (\S\ref{sec:gauge_freedom}) is the paper's
most MechRef-relevant finding. Ten seeds of modular multiplication all
achieve $\IIA$ 0.94--1.00, establishing behavioral and role-level transport:
the term ``the Grassmannian causal variable for multiplication'' refers to
the same functional role across seeds. But the recovered DAS subspaces are
nearly orthogonal (pairwise overlap $\approx 0.008$, below the random
baseline of $k/d = 0.016$), so subspace-level transport fails. Object-level
transport (same specific $Q$ matrix) fails \emph{a fortiori}. Structural
transport---whether the gauge-invariant composites ($W_Q W_K^\top$, SVD
singular vectors) are preserved---is partially supported by the SVD result
($\IIA = 0.999$ from weight structure alone) but not systematically tested
across seeds.

The implication is precise: ``the causal variable for multiplication'' refers
at the role level but not at the subspace level. This is not a negative
result---it is informative about the correct level of description. The term
refers, but it refers to a functional equivalence class, not to a geometric
object.

\paragraph{Grokking: reference is conditional.}
The stochastic class (\S\ref{sec:stochastic}) is a natural experiment in
reference failure. ``The causal variable for power'' does not refer in
general: at seed 137, it refers to a grokked representation with
equivariance 96.5\%; at seed 42, it refers to a memorized lookup table with
equivariance 66.5\%. The term's reference is conditional on grokking having
occurred, which is itself stochastic for operations near the partition
boundary. This makes the stochastic class a boundary case for MechRef:
behavioral-level transport fails across seeds because the behavioral outcome
itself is seed-dependent.

Across architectures, grokking transports behaviorally (DMF+ReLU and the
full transformer both generalize on multiplication) but only partially at
the role level: the linearized transformer groks without Fourier structure
(\S\ref{sec:dmf}), meaning the algorithm implementing generalization differs
across architectures. ``Grokking'' refers to the same training-dynamics
phenomenon (memorize-then-generalize) but not necessarily to the same
internal mechanism.

\paragraph{Structured pi-SAE variable: method-level reference.}
The Structured pi-SAE is applied to both grokking models and GPT-2, and
achieves high IIA in both settings. However, the causal variables it
recovers are not the same kind of entity: in grokking, the latent captures
Fourier-structured circular representations; in GPT-2, it captures
task-relevant features (gender, hypernymy class, numerical magnitude). The
term transports at the behavioral level (both achieve interchange success)
and the role level (both identify a task-relevant latent factor), but not at
the subspace or structural level. The method is portable; the recovered
variable is task-specific.

\paragraph{Equivariance: incommensurable across domains.}
Equivariance testing is the paper's strongest structural diagnostic for
grokking models, where the operation's symmetry group ($\Z_p$ acting on
residues) provides a natural test. But NLP tasks lack algebraic symmetries,
so the paper substitutes cross-distribution generalization (disjoint names,
novel templates) as the analogous test. These are not the same kind of
evidence: equivariance tests algebraic coherence, while cross-distribution
generalization tests distributional robustness. The term ``structural
diagnostic'' transports across domains, but the underlying measurement does
not. This is an honest limitation: the paper's strongest evidence type is
available only for the synthetic setting.
